
\section{期望清单}

本章自始至终在一个域 $k$ 上工作；$X$ 将是一个\trProjective{} $k$-\trScheme，而 $n=\dim X$。我们将会想要 $X$ 上的一个\trCoherentSheaf{} $\boxed{\omega}$（或者，说得更精确些，$\boxed{\omega_X}$；甚至，更好的是 $\omega_{X/k}$），也就是\textbf{\trCanonical{}层}或\textbf{\trDualizingSheaf}，它将在\trDuality{}的诸陈述中扮演角色。为了得到最好的那些陈述，我们还会想要一个\textbf{迹\trMorphism}
\begin{equation*}\tag{29.1.0.1}\label{eq:trace-morphism-hoped}
\boxed{t\colon\mathrm{H}^n(X,\omega_X)\longrightarrow k.}
\end{equation*}
我们最终将在 \S29.1.5 中定义 $\omega$ 与 $t$，但首先我们来讨论一下，我们希望它们具有哪些性质。

\textbf{29.1.1. \trDesideratum{}：对光滑簇而言，\trCotangentBundle{}的\trDeterminant{}是\trDualizingSheaf。}\footnote{译注：节题 Desiderata（期望清单）与本条目头 Desideratum（愿望）为同源词：本节罗列的正是我们期望对偶性应当具有的一系列性质，故条目均以「愿望」为头。——译者注}\quad 若 $X$ \trSmooth，则在此情形我们将会想要 $\boxed{\omega_X=\mathscr{K}_X}$（回想 $\mathscr{K}_X=\det\Omega_X$）——亦即「\trCanonical{}丛是 \trSerre{}对偶化的」这一奇迹（\S18.5.2）。特别地，$\omega_X$ 是一个\trInvertibleSheaf。这件事证明起来会令人不安地比我们所展示的那些基本\trDuality{}陈述更难；我们只能留待以后再处理它（\mbox{练习~29.4.K}）。但我们会证明更多的结果，比如：若 $X$ 是一个\trSmooth\trVariety{}中的\trRegularEmbedding，则 $\omega_X$ 是一个\trInvertibleSheaf（\mbox{练习~29.4.J}）。

\textbf{29.1.2. \trDesideratum{}：\trVectorBundle{}的\trSerreDuality。}\quad \trDuality{}的第一个版本——它与\mbox{\trDesideratum{}~29.1.1} 一起给出\mbox{定理~18.5.1}——如下：若 $\mathscr{F}$ 是有限秩的局部自由层，则我们有函子性的\trIsomorphism
\begin{equation*}\tag{29.1.2.1}\label{eq:vector-bundle-duality-iso}
\mathrm{H}^i(X,\mathscr{F})\cong\mathrm{H}^{n-i}(X,\mathscr{F}^{\vee}\otimes\omega_X)^{\vee}.
\end{equation*}
更精确地说，我们要构造的是\emph{某个特定的}\trIsomorphism{} \eqref{eq:vector-bundle-duality-iso}，或者等价地，某个特定的\trPerfectPairing{} $\mathrm{H}^i(X,\mathscr{F})\times\mathrm{H}^{n-i}(X,\mathscr{F}^{\vee}\otimes\omega_X)\to k$。这个\trIsomorphism{}将是随 $\mathscr{F}$ 函子的，也就是说，它给出协变函子之间的一个\trNaturalIsomorphism
\begin{equation*}\tag{29.1.2.2}\label{eq:vector-bundle-functorial-iso}
\boxed{\mathrm{H}^i(X,\cdot)\overset{\sim}{\longrightarrow}\mathrm{H}^{n-i}(X,\cdot^{\vee}\otimes\omega_X)^{\vee}.}
\end{equation*}
我们称之为\textbf{\trVectorBundle{}的函子版本\trSerreDuality}。

更好的是，\trCohomology{}中还应当有一个上积（cup product），可用它来构造一个映射 $\mathrm{H}^i(X,\mathscr{F})\times\mathrm{H}^{n-i}(X,\mathscr{F}^{\vee}\otimes\omega_X)\to\mathrm{H}^n(X,\omega_X)$。它应当随 $\mathscr{F}$ 函子（其含义是：我们得到函子之间的一个\trNaturalTransformation{} $\mathrm{H}^i(X,\cdot)\to\mathrm{H}^{n-i}(X,\cdot^{\vee}\otimes\omega_X)^{\vee}\otimes\mathrm{H}^n(X,\omega_X)$）。将它与\trTraceMap \eqref{eq:trace-morphism-hoped} 相结合，我们就得到一个映射
\begin{equation*}\tag{29.1.2.3}\label{eq:vector-bundle-trace-composition}
\begin{tikzcd}[row sep=3.2em, column sep=3.4em]
\mathrm{H}^i(X,\mathscr{F})\times\mathrm{H}^{n-i}(X,\mathscr{F}^{\vee}\otimes\omega_X) \ar[r] & \mathrm{H}^n(X,\mathscr{F}\otimes\mathscr{F}^{\vee}\otimes\omega_X) \ar[d] \\
& \mathrm{H}^n(X,\omega_X) \ar[r, "t"] & k,
\end{tikzcd}
\end{equation*}
它若是一个\trPerfectPairing，就应当给出 \eqref{eq:vector-bundle-functorial-iso}。若复合映射 \eqref{eq:vector-bundle-trace-composition} 是一个\trPerfectPairing，我们就说 $X$（连同附加数据 $(\omega_X,t)$）满足\textbf{\trVectorBundle{}\trSerreDuality{}的迹版本}。（在\mbox{定理~18.5.1} 的陈述之后，我们曾暗示过这一点。）

事实上，\trVectorBundle{}\trSerreDuality{}的迹版本（从而函子版本）在 $X$ 为 \trCM{}时成立，特别地，在 $X$ \trSmooth{}时成立。我们将证明函子版本（见\mbox{推论~29.3.10}）；至于迹版本可以如何证明，我们会在\mbox{注记~29.3.15} 中给出一点提示。

\textbf{29.1.3. \trDesideratum{}：更一般的 $X$ 与 $\mathscr{F}$ 的\trDuality。}\quad 一种较弱的\trDuality{}将在较弱的假设下成立。我们将证明：对纯 $n$ 维\trProjective{} $k$-\trScheme{} $X$ 上的任意\trCoherentSheaf{} $\mathscr{F}$（无需 \trCM{}假设），存在一个\emph{函子性的\trIsomorphism}（随 $\mathscr{F}$ 函子）
\begin{equation*}\tag{29.1.3.1}\label{eq:hom-functorial-iso}
\boxed{\mathrm{Hom}(\mathscr{F},\omega_X)\overset{\sim}{\longrightarrow}\mathrm{H}^n(X,\mathscr{F})^{\vee}.}
\end{equation*}
我们称之为\textbf{$\mathrm{Hom}$ 的函子版本\trSerreDuality}。

与\trVectorBundle{}的\trSerreDuality{}相平行，对于这个\trPerfectPairing，我们有一个自然的候选者：
\begin{equation*}\tag{29.1.3.2}\label{eq:hom-trace-pairing}
\boxed{\mathrm{Hom}(\mathscr{F},\omega_X)\times\mathrm{H}^n(X,\mathscr{F})\longrightarrow\mathrm{H}^n(X,\omega_X)\overset{t}{\longrightarrow}k,}
\end{equation*}
其中\trTraceMap $t$ 是 $\mathrm{H}^n(X,\omega_X)$ 上的某个线性泛函。若这个复合是一个\trPerfectPairing，我们就说 $X$（连同附加数据 $(\omega_X,t)$）满足\textbf{$\mathrm{Hom}$ 的\trSerreDuality{}的迹版本}。与\trVectorBundle{}的\trSerreDuality{}的迹版本不同，这里的「上积」映射 $\mathrm{Hom}(\mathscr{F},\omega_X)\times\mathrm{H}^n(X,\mathscr{F})\to\mathrm{H}^n(X,\omega_X)$ 我们已经知道是什么：$\mathrm{Hom}(\mathscr{F},\omega_X)$ 的一个元素 $[\sigma\colon\mathscr{F}\to\omega_X]$，借 $\mathrm{H}^n(X,\cdot)$ 的协变性（见 \S18.1），诱导出一个映射 $\mathrm{H}^n(X,\mathscr{F})\to\mathrm{H}^n(X,\omega_X)$。所得的配对显然随 $\mathscr{F}$ 函子，因此 $\mathrm{Hom}$ 的\trSerreDuality{}的迹版本蕴含函子版本。而与\trVectorBundle{}的情形不同的是，$\mathrm{Hom}$ 的函子版本\trSerreDuality{}反过来也蕴含迹版本：

\textbf{29.1.A.（练习）}\quad 证明：$\mathrm{Hom}$ 的\trSerreDuality{}的函子版本蕴含迹版本。换言之，\trTraceMap{}已经隐含在函子性的\trIsomorphism{} $\mathrm{Hom}(\mathscr{F},\omega_X)\overset{\sim}{\longrightarrow}\mathrm{H}^n(X,\mathscr{F})^{\vee}$ 之中。提示：考虑（由函子性得来的）交换图
\[
\begin{tikzcd}[row sep=3.2em, column sep=4em]
\mathrm{Hom}(\mathscr{F},\omega_X) \ar[r] & \mathrm{H}^n(X,\mathscr{F})^{\vee} \\
\mathrm{Hom}(\omega_X,\omega_X) \ar[u] \ar[r] & \mathrm{H}^n(X,\omega_X)^{\vee}. \ar[u]
\end{tikzcd}
\]

\textbf{29.1.4. 开始动手。}\quad 我们先来定义其中的若干术语。

\textbf{29.1.5.}\ \emph{\trhDefinition{}。}\quad 设 $X$ 是\trDimension{}为 $n$ 的\trProjective{} $k$-\trScheme。一个\trCoherentSheaf{} $\omega=\omega_X=\omega_{X/k}$ 连同一个映射 $t\colon\mathrm{H}^n(X,\omega_X)\to k$ 合起来称作\textbf{对偶化的}，若自然映射（对照 \eqref{eq:hom-functorial-iso}）
\begin{equation*}\tag{29.1.5.1}\label{eq:dualizing-pairing}
\mathrm{Hom}(\mathscr{F},\omega_X)\times\mathrm{H}^n(X,\mathscr{F})\longrightarrow\mathrm{H}^n(X,\omega_X)\overset{t}{\longrightarrow}k
\end{equation*}
对 $X$ 上的一切\trCoherentSheaf{} $\mathscr{F}$ 都是\trPerfectPairing。我们称 $\omega_X$ 为\textbf{\trCanonical{}层}，称 $t$ 为\textbf{\trTraceMap}。（早先关于 $\omega_X$ 与 $t$ 的讨论只是愿望式的；现在这则是一个定义。）若 $X$ 具有这样一个 $(\omega_X,t)$，我们就说 $X$ 满足\textbf{$\mathrm{Hom}$ 的\trSerreDuality}。

\textbf{29.1.6.}\ \emph{记号旁注。}\quad 此处定义的\trCanonical{}层常被称作\textbf{\trDualizingSheaf}。这个术语并不算糟，但最好的做法或许是：把这个称呼留给满足\trSerreDuality{}最好版本——即 $\mathrm{Ext}$ 的迹版本（\S29.1.13）——的层，而那已超出本书的范围。类似地，为了精确起见，\mbox{定义~29.1.5} 中的「对偶化的」或许更宜称作「弱对偶化的」。

此处的「\trCanonical{}层」$\omega_X$，与 \S21.5.3 中光滑簇的「\trCanonical{}丛」$\mathscr{K}_X=\det\Omega_X$，在术语上存在一个表面上的冲突。但 \S29.4 的要点正是要表明：在两者都适用的场合，这两个概念是同一个东西。

\textbf{29.1.7.}\quad 下述\trhProposition{}为「the canonical sheaf（这个\trCanonical{}层）」这一说法中使用定冠词「the」（而非「a」）提供了正当性。

\textbf{29.1.8. \trhProposition。}\ ——\ \emph{若 $X$ 存在\trCanonical{}层与迹 $(\omega_X,t)$，则这份数据在相差唯一\trIsomorphism{}的意义下唯一。}

\emph{证明。}\quad 设我们有两个这样的 $(\omega_X,t)$ 与 $(\omega'_X,t')$。由那两个\trMorphism
\begin{equation*}\tag{29.1.8.1}\label{eq:two-trace-data}
\begin{gathered}
\mathrm{Hom}(\mathscr{F},\omega_X)\times\mathrm{H}^n(X,\mathscr{F})\longrightarrow\mathrm{H}^n(X,\omega_X)\overset{t}{\longrightarrow}k,\\[1.4em]
\mathrm{Hom}(\mathscr{F},\omega'_X)\times\mathrm{H}^n(X,\mathscr{F})\longrightarrow\mathrm{H}^n(X,\omega'_X)\overset{t'}{\longrightarrow}k,
\end{gathered}
\end{equation*}
我们得到一个自然的\trBijection{} $\mathrm{Hom}(\mathscr{F},\omega_X)\cong\mathrm{Hom}(\mathscr{F},\omega'_X)$，它随 $\mathscr{F}$ 函子。由典型的\trUniversalProperty{}论证（\mbox{练习~1.2.Z}），这诱导出一个（唯一的）\trIsomorphism{} $\omega_X\overset{\sim}{\hookrightarrow}\omega'_X$。由 \eqref{eq:two-trace-data}，在这一同构之下，两个\trTraceMap $t$ 与 $t'$ 也必定相同。\hfill$\square$

我们将在\mbox{推论~29.3.12} 中证明 $\mathrm{Hom}$ 的\trSerreDuality{}的函子版本与迹版本。射影空间这一特殊情形将是关键的原料。我们现在就来证明这个情形。

\textbf{29.1.9.}\ \emph{$\mathbb{P}^n_k$ 上（$\mathrm{Hom}$ 的）\trSerreDuality。}\quad 定义 $\omega_{\mathbb{P}^n_k}$（为方便起见，或径记作 $\omega$）为 $\OO_{\mathbb{P}^n_k}(-n-1)$。令 $t$ 为任一\trIsomorphism{} $\mathrm{H}^n(\mathbb{P}^n_k,\omega_{\mathbb{P}^n_k})\overset{\sim}{\hookrightarrow}k$（\mbox{定理~18.1.2}(b)）。正如这个记号所暗示的，$(\omega_{\mathbb{P}^n_k},t)$ 对射影空间 $\mathbb{P}^n_k$ 而言将会是对偶化的。

\textbf{29.1.B.（练习）}\quad 设 $\mathscr{F}=\OO_{\mathbb{P}^n_k}(m)$。证明：自然映射 \eqref{eq:dualizing-pairing} 是一个\trPerfectPairing。（提示：亲自动手做。参看\mbox{定理~18.1.2} 之后的讨论。）进而证明：若 $\mathscr{F}$ 是 $\mathbb{P}^n_k$ 上若干\trLineBundle{}的一个\trDirectSum，则自然映射 \eqref{eq:dualizing-pairing} 是一个\trPerfectPairing。

\textbf{29.1.10. \trhProposition。}\ ——\ \emph{$\mathrm{Hom}$ 的\trSerreDuality{}的函子版本（从而由\mbox{练习~29.1.A}，迹版本亦然）对 $\mathbb{P}^n_k$ 成立。}

\emph{证明。}\quad 我们希望证明：由上积配对与\trIsomorphism{} $t$ 的选取而来的（函子性的）自然映射
\[
\mathrm{Hom}(\mathscr{F},\omega_{\mathbb{P}^n_k})\overset{\sim}{\longrightarrow}\mathrm{H}^n(\mathbb{P}^n_k,\mathscr{F})^{\vee},
\]
是一个\trIsomorphism（对照 \eqref{eq:hom-functorial-iso}）。取定 $\mathbb{P}^n_k$ 上的一个\trCoherentSheaf{} $\mathscr{F}$。依据\mbox{定理~16.1.1}，我们可以将 $\mathscr{F}$ 呈现为
\begin{equation*}\tag{29.1.10.1}\label{eq:presentation-of-f}
0\longrightarrow\mathscr{G}\longrightarrow\mathscr{E}\longrightarrow\mathscr{F}\longrightarrow 0,
\end{equation*}
其中 $\mathscr{E}$ 是有限多个\trLineBundle{}的一个\trDirectSum，而 $\mathscr{G}$ 凝聚。对 \eqref{eq:presentation-of-f} 施加\trLeftExact{}函子 $\mathrm{Hom}(\cdot,\omega_{\mathbb{P}^n_k})$，我们便有\trExactSequence
\begin{equation*}\tag{29.1.10.2}\label{eq:hom-left-exact}
0\longrightarrow\mathrm{Hom}(\mathscr{F},\omega_{\mathbb{P}^n_k})\longrightarrow\mathrm{Hom}(\mathscr{E},\omega_{\mathbb{P}^n_k})\longrightarrow\mathrm{Hom}(\mathscr{G},\omega_{\mathbb{P}^n_k}).
\end{equation*}
对 \eqref{eq:presentation-of-f} 取\trCohomology{}的\trLongExactSequence{}并对偶化，我们便有\trExactSequence
\begin{equation*}\tag{29.1.10.3}\label{eq:cohom-dual-exact}
0\longrightarrow\mathrm{H}^n(\mathbb{P}^n_k,\mathscr{F})^{\vee}\longrightarrow\mathrm{H}^n(\mathbb{P}^n_k,\mathscr{E})^{\vee}\longrightarrow\mathrm{H}^n(\mathbb{P}^n_k,\mathscr{G})^{\vee}.
\end{equation*}
（函子性的）配对 \eqref{eq:hom-trace-pairing} 给出从 \eqref{eq:hom-left-exact} 到 \eqref{eq:cohom-dual-exact} 的一个映射：
\begin{equation*}\tag{29.1.10.4}\label{eq:five-lemma-map}
\begin{tikzcd}[row sep=3.2em, column sep=2.6em]
0 \ar[r] & 0 \ar[r] & \mathrm{H}^n(\mathbb{P}^n_k,\mathscr{F})^{\vee} \ar[r] & \mathrm{H}^n(\mathbb{P}^n_k,\mathscr{E})^{\vee} \ar[r] & \mathrm{H}^n(\mathbb{P}^n_k,\mathscr{G})^{\vee} \\
0 \ar[r] \ar[u, "\alpha"] & 0 \ar[r] \ar[u, "\beta"] & \mathrm{Hom}(\mathscr{F},\omega_{\mathbb{P}^n_k}) \ar[r] \ar[u, "\gamma"] & \mathrm{Hom}(\mathscr{E},\omega_{\mathbb{P}^n_k}) \ar[r] \ar[u, "\delta"] & \mathrm{Hom}(\mathscr{G},\omega_{\mathbb{P}^n_k}). \ar[u, "\epsilon"]
\end{tikzcd}
\end{equation*}
映射 $\alpha$ 与 $\beta$ 显然是\trIsomorphism，而\mbox{练习~29.1.B} 表明 $\delta$ 是一个\trIsomorphism。于是，由五引理的微妙版本（\mbox{练习~1.6.D}；注意到 $\beta$ 与 $\delta$ 均\trInjectiveMap{}而 $\alpha$ \trSurjectiveMap），$\gamma$ \trInjectiveMap。这表明：对\emph{一切}\trCoherentSheaf{} $\mathscr{F}'$，自然映射 $\mathrm{Hom}(\mathscr{F}',\omega_{\mathbb{P}^n_k})\to\mathrm{H}^n(\mathbb{P}^n_k,\mathscr{F}')^{\vee}$ 都\trInjectiveMap；特别地，取 $\mathscr{F}'=\mathscr{G}$ 时也如此。于是 $\epsilon$ \trInjectiveMap。再由五引理微妙版本的对偶命题（由于 $\beta$ 与 $\delta$ 均\trSurjectiveMap，而 $\epsilon$ \trInjectiveMap），$\gamma$ \trSurjectiveMap。\hfill$\square$

\textbf{29.1.11.}\ \emph{数学谜题。}\quad 这里有一个谜题，逼你直面一个可能令人困惑的点。我们将会看到，\mbox{\trDesideratum{}~29.1.1} 对 $\mathbb{P}^1$ 成立，于是 $\omega_{\mathbb{P}^1}\cong\Omega_{\mathbb{P}^1}$。那么\trTraceMap $t\colon\mathrm{H}^1(\mathbb{P}^1,\Omega_{\mathbb{P}^1})\to k$ 究竟是什么？$\mathrm{H}^1(\mathbb{P}^1,\Omega_{\mathbb{P}^1})$ 的\trCech{}复形（取 $\mathbb{P}^1$ 的惯用\trCover）由下式给出
\begin{equation*}\tag{29.1.11.1}\label{eq:cech-complex-omega-p1}
0\longrightarrow\Omega_{\mathbb{P}^1}(U_0)\times\Omega_{\mathbb{P}^1}(U_1)\overset{\alpha}{\longrightarrow}\Omega_{\mathbb{P}^1}(U_0\cap U_1)\longrightarrow 0.
\end{equation*}
若 $U_0=\Spec k[x]$，且 $U_0\cap U_1=\Spec k[x,1/x]$，则 $U_0\cap U_1$ 上的那些\trDifferential{}都形如 $f(x)\,\mathrm{d}x$，其中 $f(x)$ 是一个\trLaurentPolynomial（例如：$(x^{-3}+x^{-1}+3+17x^4)\,\mathrm{d}x$）。要计算 $\mathrm{H}^1(\mathbb{P}^1,\Omega_{\mathbb{P}^1})$，我们需要找出 $U_0\cap U_1$ 上哪些这样的\trDifferential{}落在 \eqref{eq:cech-complex-omega-p1} 中 $\alpha$ 的\trImage{}里。显然，形如 $x^i\,\mathrm{d}x$（$i\ge 0$）的任何一项都能延拓为 $U_0$ 上的一个\trDifferential，从而落在 $\alpha$ 的\trImage{}中。短短一算便可知道，形如 $x^i\,\mathrm{d}x$（$i<-1$）的任何一项都能延拓为 $U_1$ 上的一个\trDifferential。于是 $\alpha$ 的\trCokernel{}可以描述为由 $x^{-1}\,\mathrm{d}x$ 生成的那个一维 $k$-\trVectorSpace。我们有一个到 $k$ 的显然\trIsomorphism：取 $x^{-1}\,\mathrm{d}x$ 的系数，这可以被解释为「取 $0$ 处的留数」。但还有另一个同等好的选择：取 $\infty$ 处的留数——当然，没有任何理由偏爱 $0$ 胜过 $\infty$（或者偏爱 $U_0$ 胜过 $U_1$）！然而这两个留数并\emph{不}相同——它们之和为 $0$（你可以很快算出来——你也可以凭\trRiemannSurface{}理论中的留数定理相信这一点）。那么：哪一个才是迹？

\textbf{29.1.12. \trDesideratum{}：\trDuality{}的一个更强的版本，涉及 $\mathrm{Ext}$。}

\epigraph{你的公式 $\mathrm{H}^{n-p}(X,\mathscr{F})'=\mathrm{Ext}^p_{\OO}(X,F,\Omega^n)$ 令我兴奋不已，因为我深信：无论在解析情形还是代数情形（后者更重要——对我而言），它都是陈述对偶定理的正确方式。}{——J.-P.~\trSerre，致 A.~\trGrothendieck{}书信，1955 年 12 月 22 日 [GrS, p.~21]}

\trVectorBundle{}与 $\mathrm{Hom}$ 的\trSerreDuality{}有一个共同的扩张。为此，我们需要 $\OO$-\trModule{}的 $\mathrm{Ext}$ 概念。我们将在 \S29.2 中发展所需的内容，但眼下只需注意到：若 $X$ 是一个\trRingedSpace，则 $\OO_X$-\trModule{}的\trCategory{}有足够多的\trInjective{}对象（\mbox{定理~23.4.1}），于是我们可以定义 $\mathrm{Ext}^i(\mathscr{F},\cdot)$ 为 $\mathrm{Hom}(\mathscr{F},\cdot)$ 的\trDerivedFunctor（正如 \S23.2.4 中定义的 $A$-\trModule{}的 $\mathrm{Ext}$-函子那样）。

若我们有函子之间的一个\trIsomorphism
\begin{equation*}\tag{29.1.12.1}\label{eq:ext-functorial-iso}
\boxed{\mathrm{Ext}^i(\cdot,\omega_X)\overset{\sim}{\longrightarrow}\mathrm{H}^{n-i}(X,\cdot)^{\vee},}
\end{equation*}
我们就说 $X$ 满足\textbf{$\mathrm{Ext}$ 的函子版本\trSerreDuality}。（$i=0$ 的情形就是 $\mathrm{Hom}$ 的函子版本\trSerreDuality，或者由\mbox{练习~29.1.A}，即迹版本。）

在\mbox{练习~29.2.K} 中，我们将发现：$\mathrm{Ext}$ 的函子版本\trSerreDuality（迹版本亦然）蕴含\trVectorBundle{}的函子版本\trSerreDuality（迹版本亦然）。于是，要证明\mbox{定理~18.5.1}，只须证明 $\mathrm{Ext}$ 的函子版本\trSerreDuality，再加上\mbox{\trDesideratum{}~29.1.1}（见\mbox{注记~29.4.10}）。当 $X$ 为 \trCM{}时，$\mathrm{Ext}$ 的函子版本\trSerreDuality{}将在\mbox{推论~29.3.14} 中得到证明。

\textbf{29.1.13.}\ \emph{$\mathrm{Ext}$ 的\trSerreDuality{}的迹版本。}\quad 与前几种版本一样，$\mathrm{Ext}$ 的函子版本\trSerreDuality{}的函子性也应当有其来源。存在上积
\[
\mathrm{Ext}^a(\mathscr{F},\mathscr{G})\times\mathrm{Ext}^b(\mathscr{E},\mathscr{F})\longrightarrow\mathrm{Ext}^{a+b}(\mathscr{E},\mathscr{F})
\]
（即「\trYoneda{}上积」）。若复合
\[
\resizebox{0.85\textwidth}{!}{\begin{tikzcd}[row sep=3.4em, column sep=1.6em, ampersand replacement=\&]
\mathrm{Ext}^i_X(\mathscr{F},\omega_X)\times\mathrm{Ext}^{n-i}_X(\mathscr{O}_X,\mathscr{F})=\mathrm{Ext}^i_X(\mathscr{F},\omega_X)\times\mathrm{Ext}^{n-i}_X(\mathscr{O}_X,\mathscr{F}) \ar[dr, "\cup"] \& \\
\& \mathrm{Ext}^n_X(\mathscr{O}_X,\omega_X)=\mathrm{H}^n(X,\omega_X) \ar[r, "t"] \& k
\end{tikzcd}}
\]
对一切 $i$ 都是\trPerfectPairing，我们就说 $X$ 满足\textbf{$\mathrm{Ext}$ 的\trSerreDuality{}的迹版本}。我们将无力证明迹版本，因为我们并不会定义这个上积。但一切纯维的\trProper{} \trCM{} $k$-\trScheme{}都满足这一强形式的\trDuality。[Stacks, tag~0FVU]

\textbf{29.1.14. \trCM{}假设的必要性。}\quad 我们指出：凡是在陈述了 \trCM{}假设的地方，这些假设都是必要的。下面的例子在一切情形都适用。设 $X$ 是 $\mathbb{P}^4_k$ 中交于一点 $p$ 的两个二维平面的并。若 $X$ 上存在一个凝聚的\trCanonical{}层 $\omega_X$，则当 $d\gg 0$ 时，由\trSerreDuality{}将会得到 $\mathrm{h}^1(X,\OO_X(-d))=\mathrm{h}^1(X,\omega_X(d))$，而由\trSerre{}消失（\mbox{定理~18.1.3}(ii)）它必须为 $0$。我们将证明事实并非如此。

\textbf{29.1.C.（练习）}\quad 设 $H$ 是 $\mathbb{P}^4$ 中不经过 $p$ 的一个\trHyperplane；比如说，$H$ 就是超平面 $x_0=0$。设 $Z$ 是 $dH$（除子 $x_0^d=0$）与 $X$ 的交。\footnote{译注：原书此处作「for \(d\ge 0\)」；然 $d=0$ 时 $Z=X$ 而 $\mathrm{h}^0(Z,\OO_Z)=1$，疑当为「d > 0」，兹照录原书。——译者注}（用廉价的方式）证明：当 $d\ge 0$ 时，$\mathrm{h}^0(Z,\OO_Z)>1$。利用\trExactSequence
\[
0\longrightarrow\OO_X(-Z)\longrightarrow\OO_X\longrightarrow\OO_Z\longrightarrow 0
\]
来证明 $\mathrm{h}^1(X,\OO_X(-d))>0$。
